\chapter{Li Negative Inversion and the Global Weil Frontier}
\label{chap:li-negative-inverse}

\section{Canonical scope of the promoted Li/Weil line}

This chapter is the integration point for the Version 0.9 promoted claim line SOH-L012--SOH-L032. Earlier development snapshots used some identifiers in that numerical interval for Hermite, reciprocal-tail, and DHSE results. Those historical meanings are preserved as provenance but are migrated to domain identifiers in \cref{app:claims}; they do not redefine the promoted identifiers used here.

The purpose of the integration is twofold. First, it places the critical-axis geometry, the negative inverse, Li's criterion, and Weil positivity in one coordinate chain. Second, it makes the remaining proof obligation in the Li/Weil branch explicit. No result in this chapter promotes SOH-C005 or the Riemann Hypothesis to proved status. Later chapters develop a separate quotient/PF branch with its own RH-equivalent frontier SOH-G003.

\section{The critical-axis projective coordinate}

Define
\begin{equation}
 \Omega(s)=\frac{s}{1-s},
 \qquad
 K(s)=1-\overline{s}.
 \label{eq:v2-omega-K}
\end{equation}
A direct calculation gives
\begin{equation}
 \Omega(K(s))=\frac{1}{\overline{\Omega(s)}}.
 \label{eq:v2-omega-K-conjugacy}
\end{equation}
For $s=\sigma+it$,
\begin{equation}
 |\Omega(s)|^2
 =\frac{\sigma^2+t^2}{(1-\sigma)^2+t^2}.
 \label{eq:v2-omega-modulus}
\end{equation}
Hence
\begin{equation}
 \boxed{\Re s=\frac12\iff |\Omega(s)|=1.}
 \label{eq:v2-critical-unit-circle}
\end{equation}
Thus the conjugate-affine, anti-holomorphic involution $K$ becomes reciprocal-conjugate inversion in $\Omega$-space, and its fixed locus becomes the unit circle. In the centered variable $s-1/2$ the same action is genuinely anti-linear.

Define the signed log-radial field and its non-negative potential
\begin{equation}
 B(s)=\log|\Omega(s)|,
 \qquad
 V(s)=B(s)^2.
 \label{eq:v2-BV}
\end{equation}
Then
\begin{equation}
 B(K(s))=-B(s),
 \qquad
 V(K(s))=V(s),
 \qquad
 V(s)=0\iff\Re s=\frac12.
 \label{eq:v2-BV-symmetry}
\end{equation}
These identities are exact coordinate statements. In particular, $V$ is not being introduced as an independent physical potential.

\section{The negative inverse is the Li coordinate}

Li's zero coordinate is
\begin{equation}
 z_L(s)=1-\frac1s=\frac{s-1}{s}.
 \label{eq:v2-li-coordinate}
\end{equation}
Combining this with \cref{eq:v2-omega-K} gives the exact identity
\begin{equation}
 \boxed{z_L(s)=-\frac1{\Omega(s)}.}
 \label{eq:v2-negative-inverse}
\end{equation}
Thus the negative inverse is not an additional fitted map: it is exactly Li's coordinate after the projective critical-axis transform. Moreover,
\begin{equation}
 z_L(K(s))=\frac1{\overline{z_L(s)}},
 \qquad
 \Re s=\frac12\iff|z_L(s)|=1,
 \label{eq:v2-li-K}
\end{equation}
and
\begin{equation}
 B(s)=-\log|z_L(s)|,
 \qquad
 V(s)=\log^2|z_L(s)|.
 \label{eq:v2-li-BV}
\end{equation}

\section{Reciprocal-conjugate quartets and local Li growth}

With the standard symmetric zero-sum convention, Li's coefficients may be written \citep{li1997}
\begin{equation}
 \lambda_n
 =\sum_{\rho}\left[1-\left(1-\frac1\rho\right)^n\right]
 =\sum_{\rho}[1-z_L(\rho)^n].
 \label{eq:v2-li-coefficients}
\end{equation}
For one reciprocal-conjugate quartet
\begin{equation}
 Q(z)=\{z,\overline z,z^{-1},\overline z^{-1}\},
 \qquad z=Re^{i\phi},\quad R>0,
 \label{eq:v2-quartet}
\end{equation}
the local contribution is exactly
\begin{equation}
 \boxed{
 L_n(Q)=4-2(R^n+R^{-n})\cos(n\phi).
 }
 \label{eq:v2-local-li}
\end{equation}
Define
\begin{equation}
 \mathscr R(\rho)=
 \limsup_{n\to\infty}
 \left(\frac{|L_n(Q)-4|}{2}\right)^{1/n}.
 \label{eq:v2-growth-radius-def}
\end{equation}
The elementary root-growth calculation gives
\begin{equation}
 \boxed{\mathscr R(\rho)=\max(R,R^{-1})=e^{|B(\rho)|}},
 \qquad
 \boxed{V(\rho)=\bigl(\log\mathscr R(\rho)\bigr)^2}.
 \label{eq:v2-growth-radius}
\end{equation}
If $R\ne1$, there are infinitely many $n$ for which $\cos(n\phi)$ is bounded away from zero on the positive side; the exponentially growing factor then forces
\begin{equation}
 \boxed{|z_L(\rho)|\ne1\Longrightarrow
 \liminf_{n\to\infty}L_n(Q)=-\infty.}
 \label{eq:v2-local-negative-subsequence}
\end{equation}
This is a local quartet theorem. It is not by itself a theorem about the sign of the global coefficient $\lambda_n$, because the global sum involves all zeros and its symmetric limiting convention.

\section{The global Li bridge}

Near $z=0$, the completed xi function has the standard Li generating form
\begin{equation}
 \log\xi\!\left(\frac1{1-z}\right)
 =\log\xi(1)+\sum_{n\ge1}\frac{\lambda_n}{n}z^n.
 \label{eq:v2-li-generating}
\end{equation}
A zero $\rho$ maps to the singular point
\begin{equation}
 z_\rho=1-\frac1\rho=z_L(\rho).
 \label{eq:v2-li-singularity}
\end{equation}
An off-axis zero with $\Re\rho>1/2$ therefore produces a singularity strictly inside the unit disk. Li's theorem gives the exact global criterion
\begin{equation}
 \boxed{
 \text{RH}\iff\lambda_n\ge0\quad\text{for every }n\ge1.
 }
 \label{eq:v2-li-criterion}
\end{equation}
No separate radius-of-convergence argument is needed here; the monograph relies on the cited Li criterion for the global equivalence rather than introducing an uncited auxiliary shortcut.

\section{Weil positivity is the same global obstruction}

Weil's explicit-formula criterion recasts the same global problem as non-negativity of an arithmetic quadratic form on an admissible test class \citep{weil1952}. In the normalized notation of the PhaseNav--Weil programme, SOH-C005 asks for an unconditional proof that
\begin{equation}
 \boxed{\mathcal W[|f|^2]\ge0
 \quad\text{for every admissible }f.}
 \label{eq:v2-weil-global-positive}
\end{equation}
Under the standard admissibility and dense-core completion, this positivity criterion is equivalent to RH. Consequently SOH-C005 is not an auxiliary lemma that can be assumed and then used to prove RH; it is the RH-hard global positivity statement expressed in the arithmetic operator language.

The Hermite ladder, fixed-section prime-tail certificate, and basis-adaptive cutoff theorem in \cref{ch:hermite-ladder,ch:prime-tail-certificate,ch:adaptive-cutoff} make the finite arithmetic side reproducible. They do not establish \cref{eq:v2-weil-global-positive} for the complete admissible family.

\section{Log-radial prime shifts and a termwise no-go lemma}

The coordinate $B=\log|\Omega|$ is additive under radial multiplication. In the prime-side screw/Weil representation the arithmetic support occurs at additive displacements $\pm\log n$. This places the prime transfer structure in the same log-radial coordinate class without identifying prime shifts with zeta zeros.

A tempting shortcut would be to argue that the positive weights $\Lambda(n)$ make every prime contribution positive semidefinite. This is false. Let
\begin{equation}
 h_a(t)=(|t|-a)_+,
 \qquad
 K_a(x,y)=h_a(x-y)-h_a(x)-h_a(y)+h_a(0).
 \label{eq:v2-hinge-kernel}
\end{equation}
Choose $x=a/2$ and $y=3a/2$. The corresponding two-point matrix is
\begin{equation}
 \begin{pmatrix}
 0&-a/2\\
 -a/2&-a
 \end{pmatrix},
 \qquad
 \det=-\frac{a^2}{4}<0.
 \label{eq:v2-hinge-indefinite}
\end{equation}
Therefore a single positive-weight hinge is indefinite. Positivity of the von Mangoldt coefficients alone cannot establish SOH-C005 term by term; any successful proof must use the complete prime, archimedean, conductor, boundary, and regularization structure.

\section{Localized arithmetic operators and the spectral floor}

For a localization scale $a>0$, write the localized arithmetic Weil form as $Q_W^a$ and define its spectral floor
\begin{equation}
 \lambda(a)=
 \inf_{0\ne f\in C_c^\infty(-a,a)}
 \frac{Q_W^a(f)}{\|f\|_2^2}.
 \label{eq:v2-spectral-floor}
\end{equation}
The localized self-adjoint realization and the Rayleigh characterization are available without assuming RH in the screw-function construction of Suzuki \citep{suzuki2026weil}. Hence the localized positivity problem is exactly
\begin{equation}
 A_a\ge0\quad\Longleftrightarrow\quad\lambda(a)\ge0.
 \label{eq:v2-local-operator-positive}
\end{equation}
For $0<a_1<a_2$, the test spaces are nested,
\begin{equation}
 C_c^\infty(-a_1,a_1)\subset C_c^\infty(-a_2,a_2),
\end{equation}
so the Rayleigh infimum obeys
\begin{equation}
 \boxed{\lambda(a_2)\le\lambda(a_1).}
 \label{eq:v2-spectral-floor-monotone}
\end{equation}
This monotonicity is a consequence of domain inclusion for the same localized Weil form; it does not supply the missing lower bound.

\section{Canonical promoted claim ledger for this branch}

The promoted statements in this chapter are SOH-L012--SOH-L032 as listed in Appendix~\ref{app:claims}. Their mathematical meanings are unchanged by the later semantic audit. The labels ``V2'' in historical receipts refer to the promotion stage at which these claims were first integrated; Version V3 retains them without reusing their identifiers.

\section{The remaining proof obligation in the Weil/Li branch}

After the reductions above, the open frontier in this branch is not another coordinate identity. It is the independent global inequality
\begin{equation}
 \boxed{
 \mathcal W[|f|^2]\ge0\quad\text{for every admissible }f,
 }
 \label{eq:v2-final-frontier}
\end{equation}
or equivalently global Li positivity for every order. A factorization such as $\mathcal W[|f|^2]=\|Af\|^2$ would close the problem only if $A$ were constructed independently from arithmetic/theta data; defining $A$ as a square root of an operator already assumed positive would be circular.

This is a branch-local statement. Chapters~26--46 develop a separate quotient/kernel/PF line whose RH-equivalent target is SOH-G003. Neither open frontier is promoted by the existence of the other.

\status{SOH-L012--SOH-L032 are retained as exact promoted reductions and no-go statements under their stated assumptions. SOH-C005 remains OPEN in the Weil/Li branch. SOH-G003 remains OPEN in the quotient branch. The Riemann Hypothesis remains OPEN.}
